\chapter{Literature Review}

This chapter reviews studies of HPV knowledge, attitudes and vaccine acceptance, with emphasis on Bangladesh and on university and young adult populations. The studies are grouped by theme, and Table~\ref{tab:lit} summarises the main ones.

\section{Awareness and Knowledge of Cervical Cancer and HPV in Bangladesh}
Awareness of cervical cancer is uneven in Bangladesh, and knowledge of HPV is generally weaker than awareness of the cancer itself. In a survey of 1,090 women at a medical college in Tangail, only 45.2\% had heard of cervical cancer, with higher awareness among urban women, those with university education and those from higher-income families~\cite{alam2022}. In a population-based survey of 2,037 ever-married women, most had heard of cervical cancer (90--93\%), but only 9\% had detailed knowledge of risk factors and 4--6\% of prevention. One in five urban women and one in twenty rural women had heard of the HPV vaccine~\cite{islam2018}. Among 600 rural women, 71.8\% were aware of cervical cancer, but knowledge of symptoms, risk factors and prevention was limited and only 2.3\% had been screened~\cite{banik2022}.

Educated women show a similar pattern. Among urban professional women working in private banks, 98\% had heard of cervical cancer but only 51\% knew that HPV is a cause, and fewer than 1\% had ever had a Pap smear~\cite{bhuiyan2018}. In a small survey of educated women, over 80\% had heard of cervical cancer but only 32.3\% had heard of the HPV vaccine~\cite{taniya2023}. Education therefore raises general awareness but does not guarantee specific knowledge of the vaccine.

\section{Knowledge and Attitudes Among University and Medical Students}
Studies of students show higher awareness but persistent gaps. A survey of 626 medical and dental students and professionals in Bangladesh found good knowledge in 43.3\% and a positive attitude to vaccination in 75.9\%, while only 11.8\% reported good practice, meaning that they had started or completed vaccination~\cite{chowdhury2022}. In Khulna, a survey of 180 female medical students reported adequate knowledge in 97.2\% and positive attitudes in 99.5\%, yet only 16.1\% had received the vaccine~\cite{promi2025}. Both studies show that knowledge and favourable attitudes do not translate directly into vaccination.

Among general female university students, a survey of 800 students at Jahangirnagar University found that students with sexual health education had higher knowledge (23.4\% good knowledge against 11.0\%), and that 64.7\% were unsure whether the vaccine protects against all HPV types that cause cervical cancer~\cite{emon2025}. In a structural equation model, attitude was the strongest predictor of willingness to be vaccinated, ahead of knowledge~\cite{emon2025}. This suggests that willingness is shaped more by how people feel about the vaccine than by factual knowledge alone.

Outside Bangladesh, a survey of 884 students at Peking University found that 72.6\% had heard of HPV and the vaccine but their knowledge was insufficient, and only 9.5\% of female students had been vaccinated. Uptake was predicted by a good relationship with doctors, sexual experience and a high level of knowledge, and limited knowledge was the main reason for refusing the vaccine~\cite{liu2020}. Most students wanted college education to be their main source of vaccine information.

\section{Willingness to Be Vaccinated}
Studies that measure willingness directly report generally high intention alongside very low uptake. Among 600 rural women, only 5.3\% had been vaccinated, but 76.6\% of the unvaccinated said they were willing; the main reasons for unwillingness were cost and concerns about the vaccine~\cite{banik2022}. In the survey of educated women, only 11\% had been vaccinated but 77\% of the rest were willing~\cite{taniya2023}. Urban professional women expressed moderate intention to accept the vaccine for themselves and high intention for their children~\cite{bhuiyan2018}. Willingness to be vaccinated is thus consistently high across different groups of women, whereas actual vaccination depends on access, cost and advice from providers.

\section{National Vaccination Programme and Remaining Gaps}
The national campaign reached about 93\% of eligible girls aged 10--14 years, helped by community engagement, work with schools and religious institutions, and responses to early rumours~\cite{dhakatribune2024,gavi2023,unbd2024}. Its design means that older adolescents and women are outside the programme, and stigma, rural access problems and poor understanding of HPV transmission continue to limit population coverage~\cite{salauddin2025}. International guidance also supports catch-up vaccination for young women up to age 26~\cite{markowitz2014}.

\section{Summary and Gap}
Three points emerge. First, awareness of HPV and the vaccine in Bangladesh is incomplete even among educated women, and detailed knowledge is poor. Second, attitudes towards vaccination are generally positive, and willingness is high among unvaccinated women, but uptake is low. Third, most evidence on students comes from medical students or focuses on knowledge. Few studies look at \emph{acceptance} among female university students from a range of academic fields, or at what distinguishes students who would accept from those who are unsure. This study addresses that gap.

\begin{longtable}{@{}>{\raggedright\arraybackslash}p{2.9cm}>{\raggedright\arraybackslash}p{2.6cm}>{\raggedright\arraybackslash}p{2.1cm}>{\raggedright\arraybackslash}p{6.0cm}@{}}
\caption{Summary of key studies reviewed.}\label{tab:lit}\\
\toprule
Study & Population & Design / $n$ & Main findings \\
\midrule
\endfirsthead
\caption[]{(continued)}\\
\toprule
Study & Population & Design / $n$ & Main findings \\
\midrule
\endhead
\bottomrule
\endlastfoot
\cite{alam2022} & Women, Tangail medical college & Cross-sectional, 1,090 & 45.2\% had heard of cervical cancer; awareness higher in urban, university-educated, higher-income women. \\
\cite{islam2018} & Ever-married women, urban and rural & Cross-sectional, 2,037 & High awareness of cervical cancer but very low detailed knowledge of risk factors (about 9\%) and prevention (4--6\%); few had heard of the vaccine. \\
\cite{banik2022} & Rural women & Cross-sectional, 600 & 71.8\% aware; 5.3\% vaccinated; 76.6\% of unvaccinated willing; cost a main barrier. \\
\cite{bhuiyan2018} & Urban professional women & Mixed methods, 160 & 98\% had heard of cervical cancer, 51\% knew HPV was a cause; 2\% vaccinated. \\
\cite{chowdhury2022} & Medical and dental students and professionals & Cross-sectional, 626 & 43.3\% good knowledge; 75.9\% positive attitude; 11.8\% good practice. \\
\cite{promi2025} & Female medical students, Khulna & Cross-sectional, 180 & High knowledge and attitude scores; 16.1\% vaccinated. \\
\cite{emon2025} & Female students, Jahangirnagar University & Cross-sectional, 800 & Sexual health education linked to higher knowledge; attitude the strongest predictor of willingness. \\
\cite{taniya2023} & Educated women & Questionnaire survey, 91 & 32.3\% had heard of the HPV vaccine; 11\% vaccinated; 77\% of unvaccinated willing. \\
\cite{liu2020} & University students, Beijing & Cross-sectional, 884 & 72.6\% had heard of HPV; 9.5\% of females vaccinated; knowledge and doctor relationship predicted uptake. \\
\end{longtable}
